\documentclass[12pt,a4paper]{article}
\usepackage[utf8]{inputenc}
\usepackage[T1]{fontenc}
\usepackage[french]{babel}
\usepackage{geometry}
\usepackage{titlesec}
\usepackage{enumitem}
\usepackage{xcolor}
\usepackage{mdframed}
\usepackage{booktabs}
\usepackage{array}
\usepackage{tabularx}
\usepackage{hyperref}
\usepackage{fontawesome5}
\usepackage{fancyhdr}
\usepackage{tikz}
\usepackage{graphicx}

\geometry{top=2.5cm, bottom=2.5cm, left=2.5cm, right=2.5cm}

% Couleurs
\definecolor{primaryBlue}{RGB}{0, 84, 166}
\definecolor{accentGreen}{RGB}{0, 153, 102}
\definecolor{warningOrange}{RGB}{230, 126, 34}
\definecolor{dangerRed}{RGB}{192, 57, 43}
\definecolor{lightGray}{RGB}{245, 245, 245}
\definecolor{darkGray}{RGB}{80, 80, 80}
\definecolor{dayColor}{RGB}{41, 128, 185}

% En-tête et pied de page
\pagestyle{fancy}
\fancyhf{}
\fancyhead[L]{\textcolor{primaryBlue}{\textbf{Système CNN — Détection de Tumeurs}}}
\fancyhead[R]{\textcolor{darkGray}{\today}}
\fancyfoot[C]{\textcolor{darkGray}{\thepage}}
\renewcommand{\headrulewidth}{1.5pt}
\renewcommand{\headrule}{\hbox to\headwidth{\color{primaryBlue}\leaders\hrule height \headrulewidth\hfill}}

% Environnement journée
\newmdenv[
  linecolor=dayColor,
  linewidth=2pt,
  topline=true,
  bottomline=true,
  leftline=true,
  rightline=false,
  backgroundcolor=lightGray,
  innerleftmargin=12pt,
  innertopmargin=8pt,
  innerbottommargin=8pt,
]{daybox}

% Environnement alerte
\newmdenv[
  linecolor=warningOrange,
  linewidth=1.5pt,
  backgroundcolor=orange!8,
  innerleftmargin=10pt,
  innertopmargin=8pt,
  innerbottommargin=8pt,
]{alertbox}

% Style des titres
\titleformat{\section}{\Large\bfseries\color{primaryBlue}}{}{0em}{}[\titlerule]
\titleformat{\subsection}{\large\bfseries\color{dayColor}}{}{0em}{\faCalendarAlt\quad}

% -------------------------------------------------------------------
\begin{document}

% Page de titre
\begin{titlepage}
  \centering
  \vspace*{2cm}
  {\Huge\bfseries\color{primaryBlue} Système de Détection et \\ Classification des Stades de Tumeurs\par}
  \vspace{0.5cm}
  {\Large\color{darkGray} Pipeline Deep Learning — Plan de Réalisation sur 4 Jours\par}
  \vspace{1.5cm}
  \begin{tikzpicture}
    \draw[primaryBlue, line width=2pt] (0,0) -- (14,0);
  \end{tikzpicture}
  \vspace{1.5cm}

  \begin{tabular}{ll}
    \textbf{Technologie principale :} & TensorFlow / Keras \\
    \textbf{Déploiement :}            & FastAPI + Docker + Streamlit \\
    \textbf{XAI :}                    & Grad-CAM \\
    \textbf{Données :}                & Kaggle Brain Tumor MRI, Breast Histopathology \\
    \textbf{Durée :}                  & 4 jours \\
  \end{tabular}

  \vfill
  {\large\color{darkGray} \today}
\end{titlepage}

% -------------------------------------------------------------------
\tableofcontents
\newpage

% -------------------------------------------------------------------
\section{Vue d'ensemble du Projet}

L'objectif est de concevoir un système robuste d'aide au diagnostic clinique basé sur le Deep Learning. Le système assure :

\begin{itemize}[leftmargin=1.5cm]
  \item \textbf{Détection binaire} : présence / absence de tumeur.
  \item \textbf{Classification multi-classes} : stades pathologiques 0 à IV (5 classes).
  \item \textbf{Explicabilité (XAI)} : cartes de chaleur Grad-CAM pour la validation médicale.
  \item \textbf{Détection d'anomalies} : auto-encodeur pour les cas atypiques.
  \item \textbf{Déploiement industrialisé} : API FastAPI conteneurisée sous Docker.
\end{itemize}

\vspace{0.5cm}
\begin{center}
\begin{tabularx}{\textwidth}{|c|X|c|}
\hline
\rowcolor{primaryBlue!20}
\textbf{Stade} & \textbf{Description Clinique} & \textbf{Focus Imagerie} \\
\hline
Stade 0  & Absence de tumeur (contrôle négatif) & — \\
\hline
Stade I  & Masse minime, asymptomatique & Micro-calcifications / nodules isolés \\
\hline
Stade II & Croissance locale, douleurs légères & Extension limitée à l'organe d'origine \\
\hline
Stade III & Envahissement local, symptômes marqués & Atteinte des ganglions lymphatiques \\
\hline
Stade IV & Métastases, signes systémiques & Lésions multi-focales, foyers secondaires \\
\hline
\end{tabularx}
\end{center}

% -------------------------------------------------------------------
\section{Stratégie des Données — Dataset et Remapping}

\subsection*{Dataset Choisi}

\begin{alertbox}
\textbf{Dataset :} \texttt{masoudnickparvar/brain-tumor-mri-dataset} (Kaggle) \\
\textbf{Lien :} \texttt{https://www.kaggle.com/datasets/masoudnickparvar/brain-tumor-mri-dataset} \\
\textbf{Volume :} $\sim$7\,023 images IRM cérébrales — $\sim$150 MB \\
\textbf{Téléchargement :} \texttt{kaggle datasets download -d masoudnickparvar/brain-tumor-mri-dataset}
\end{alertbox}

\vspace{0.4cm}

\subsection*{Problème et Solution : Remapping vers 5 Stades}

Le dataset original contient \textbf{4 dossiers/classes}. Notre projet exige \textbf{5 classes} (stades 0 à IV). La stratégie consiste à remapper les classes médicalement selon les grades WHO d'agressivité tumorale :

\vspace{0.3cm}
\begin{center}
\begin{tabular}{|l|c|l|l|}
\hline
\rowcolor{primaryBlue!20}
\textbf{Dossier Original} & \textbf{Nb Images} & \textbf{Stade Assigné} & \textbf{Justification Clinique} \\
\hline
\texttt{no\_tumor/}   & $\sim$500  & \textbf{Stade 0} & Contrôle négatif — aucune masse détectée \\
\hline
\texttt{meningioma/}  & $\sim$937  & \textbf{Stade I} & Tumeur bénigne, croissance lente (WHO Grade I) \\
\hline
\texttt{pituitary/}   & $\sim$901  & \textbf{Stade II} & Extension locale limitée, pronostic favorable \\
\hline
\texttt{glioma/} (50\% premiers) & $\sim$826 & \textbf{Stade III} & Gliome anaplasique (WHO Grade III) \\
\hline
\texttt{glioma/} (50\% derniers) & $\sim$826 & \textbf{Stade IV} & Glioblastome multiforme — GBM (WHO Grade IV) \\
\hline
\end{tabular}
\end{center}

\vspace{0.4cm}

\subsection*{Script de Remapping — \texttt{src/remap\_dataset.py}}

Le script Python suivant restructure automatiquement le dataset en 5 dossiers stades :

\begin{verbatim}
import os, shutil, glob

DATA_RAW  = "data/raw/brain-tumor-mri-dataset"
DATA_OUT  = "data/staged"
STAGES    = ["stage_0","stage_1","stage_2","stage_3","stage_4"]

mapping = {
    "no_tumor":  "stage_0",
    "meningioma":"stage_1",
    "pituitary": "stage_2",
}

for s in STAGES:
    os.makedirs(f"{DATA_OUT}/{s}", exist_ok=True)

# Stades 0, 1, 2 : copie directe
for src_class, dst_stage in mapping.items():
    for split in ["Training", "Testing"]:
        imgs = glob.glob(f"{DATA_RAW}/{split}/{src_class}/*.jpg")
        for img in imgs:
            shutil.copy(img, f"{DATA_OUT}/{dst_stage}/")

# Stade 3 et 4 : subdivision du dossier glioma (50/50)
for split in ["Training", "Testing"]:
    glioma_imgs = sorted(glob.glob(
        f"{DATA_RAW}/{split}/glioma/*.jpg"))
    mid = len(glioma_imgs) // 2
    for img in glioma_imgs[:mid]:
        shutil.copy(img, f"{DATA_OUT}/stage_3/")
    for img in glioma_imgs[mid:]:
        shutil.copy(img, f"{DATA_OUT}/stage_4/")

print("Remapping terminé.")
\end{verbatim}

\vspace{0.3cm}
\begin{center}
\begin{tabular}{|l|c|}
\hline
\rowcolor{accentGreen!20}
\textbf{Stade Final} & \textbf{Nb Images estimé} \\
\hline
\texttt{stage\_0/} (no tumor)   & $\sim$500  \\
\texttt{stage\_1/} (meningioma) & $\sim$937  \\
\texttt{stage\_2/} (pituitary)  & $\sim$901  \\
\texttt{stage\_3/} (glioma 50\%) & $\sim$826  \\
\texttt{stage\_4/} (glioma 50\%) & $\sim$826  \\
\hline
\textbf{Total}                  & \textbf{$\sim$3\,990} \\
\hline
\end{tabular}
\end{center}

\vspace{0.3cm}
\begin{alertbox}
\textbf{\faExclamationTriangle\quad Note importante :} La subdivision 50/50 du dossier \texttt{glioma} est une \textbf{approximation pédagogique} acceptable pour un projet académique. Dans un contexte clinique réel, la distinction Grade III / Grade IV nécessiterait des métadonnées radiologiques certifiées.
\end{alertbox}

% -------------------------------------------------------------------
\newpage
\section{Plan de Réalisation — 4 Jours}

% -------------------------------------------------------------------
\subsection{Jour 1 — MVP Baseline : Pipeline de Données et CNN}

\begin{daybox}
\textbf{Objectif :} Avoir un premier modèle entraîné et évalué de bout en bout.
\end{daybox}

\vspace{0.4cm}

\begin{enumerate}[label=\textbf{Étape \arabic*.}, leftmargin=2.5cm, itemsep=0.4cm]

  \item \textbf{Création de la structure du projet}

  Mettre en place l'arborescence suivante et le fichier \texttt{requirements.txt} :
  \begin{itemize}
    \item \texttt{data/} — datasets bruts et prétraités
    \item \texttt{models/} — poids sauvegardés (\texttt{.h5} / \texttt{.keras})
    \item \texttt{src/} — code Python (pipeline, modèles, entraînement)
    \item \texttt{api/} — service FastAPI
    \item \texttt{dashboard/} — application Streamlit
    \item \texttt{docker/} — Dockerfile et docker-compose
  \end{itemize}
  Dépendances principales : \texttt{tensorflow}, \texttt{opencv-python}, \texttt{scikit-learn}, \texttt{matplotlib}, \texttt{seaborn}, \texttt{kaggle}.

  \item \textbf{Pipeline de données — \texttt{src/data\_pipeline.py}}
  \begin{itemize}
    \item Exécution préalable de \texttt{src/remap\_dataset.py} pour créer les 5 dossiers stades.
    \item Chargement depuis \texttt{data/staged/} via \texttt{ImageDataGenerator.flow\_from\_directory}.
    \item Redimensionnement de toutes les images à \textbf{128×128 pixels}.
    \item Normalisation des pixels dans l'intervalle $[0, 1]$ (\texttt{rescale=1./255}).
    \item Encodage automatique des étiquettes en \textit{one-hot} sur \textbf{5 classes} (stade\_0 à stade\_4).
    \item Augmentation : rotation $\pm$20°, flip horizontal/vertical, zoom 10\%, cisaillement 10\%.
    \item Split : \textbf{70\% train / 15\% validation / 15\% test} avec \texttt{validation\_split=0.15}.
  \end{itemize}

  \item \textbf{Architecture CNN Baseline — \texttt{src/model\_baseline.py}}

  Modèle \texttt{Sequential} Keras conforme au cahier des charges :

  \vspace{0.3cm}
  \begin{center}
  \begin{tabular}{|l|l|l|}
  \hline
  \rowcolor{primaryBlue!15}
  \textbf{Couche} & \textbf{Paramètres} & \textbf{Activation} \\
  \hline
  Input       & (128, 128, 3)            & —       \\
  Conv2D      & 32 filtres, noyau 3×3    & ReLU    \\
  MaxPooling2D & Pool size (2, 2)        & —       \\
  Conv2D      & 64 filtres, noyau 3×3    & ReLU    \\
  MaxPooling2D & Pool size (2, 2)        & —       \\
  Flatten     & Vecteur 1D               & —       \\
  Dense       & 128 neurones             & ReLU    \\
  Dense (sortie) & 5 neurones            & Softmax \\
  \hline
  \end{tabular}
  \end{center}
  \vspace{0.3cm}
  Compilation : optimiseur \textbf{Adam}, perte \textbf{categorical\_crossentropy}.

  \item \textbf{Entraînement — \texttt{src/train.py}}
  \begin{itemize}
    \item Callbacks : \texttt{ModelCheckpoint} (meilleur modèle), \texttt{EarlyStopping}.
    \item Génération des courbes loss / accuracy (train vs. validation).
    \item Sauvegarde du modèle au format \texttt{.h5}.
  \end{itemize}

  \item \textbf{Évaluation — \texttt{src/evaluate.py}}
  \begin{itemize}
    \item Métriques : Accuracy, Sensibilité (Recall), F2-Score, ROC-AUC par stade.
    \item Matrice de confusion visualisée avec \texttt{seaborn}.
  \end{itemize}

\end{enumerate}

% -------------------------------------------------------------------
\newpage
\subsection{Jour 2 — XAI et Détection d'Anomalies}

\begin{daybox}
\textbf{Objectif :} Rendre le modèle explicable et capable de détecter les cas atypiques.
\end{daybox}

\vspace{0.4cm}

\begin{enumerate}[label=\textbf{Étape \arabic*.}, leftmargin=2.5cm, itemsep=0.4cm, start=6]

  \item \textbf{Grad-CAM — \texttt{src/gradcam.py}}
  \begin{itemize}
    \item Calcul des gradients de la classe cible par rapport à la dernière couche Conv2D.
    \item Génération de heatmaps pondérées et superposition sur l'image d'entrée.
    \item Utilisation clinique : mise en évidence des zones métastatiques (Stade IV), micro-calcifications (Stade I), etc.
  \end{itemize}

  \item \textbf{Auto-encodeur — \texttt{src/autoencoder.py}}
  \begin{itemize}
    \item Architecture convolutive Encoder–Decoder entraînée sur les images saines.
    \item Calcul de l'erreur de reconstruction (MSE) sur les images de test.
    \item Seuillage : erreur élevée $\Rightarrow$ cas \textbf{« Suspect / Inconnu »} $\Rightarrow$ alerte praticien.
  \end{itemize}

  \item \textbf{Incertitude Probabiliste — mise à jour de \texttt{src/evaluate.py}}
  \begin{itemize}
    \item Analyse du vecteur Softmax : si $\max(\text{softmax}) < 0.7$ $\Rightarrow$ demande de \textbf{révision manuelle prioritaire}.
    \item Journalisation des cas incertains dans un fichier de rapport.
  \end{itemize}

\end{enumerate}

% -------------------------------------------------------------------
\subsection{Jour 3 — Version Avancée : Transfer Learning et Gestion du Déséquilibre}

\begin{daybox}
\textbf{Objectif :} Dépasser les performances du CNN baseline grâce au Transfer Learning et corriger les biais liés au déséquilibre des classes.
\end{daybox}

\vspace{0.4cm}

\begin{enumerate}[label=\textbf{Étape \arabic*.}, leftmargin=2.5cm, itemsep=0.4cm, start=9]

  \item \textbf{Transfer Learning — \texttt{src/model\_advanced.py}}
  \begin{itemize}
    \item Backbone : \textbf{EfficientNetB0} ou \textbf{ResNet50V2} (poids ImageNet).
    \item \textbf{Phase 1 — Feature Extraction} : toutes les couches gelées, seule la tête de classification est entraînée.
    \item \textbf{Phase 2 — Fine-Tuning} : dégel progressif des couches supérieures pour adapter les représentations à l'imagerie médicale (niveaux de gris, contrastes spécifiques).
  \end{itemize}

  \item \textbf{Gestion du Déséquilibre — \texttt{src/losses.py}}
  \begin{itemize}
    \item Implémentation de la \textbf{Focal Loss} : $FL(p_t) = -\alpha_t (1 - p_t)^\gamma \log(p_t)$.
    \item Alternative : \textbf{Weighted Cross-Entropy} avec pondération inverse à la fréquence des classes.
    \item Objectif : pénaliser davantage les erreurs sur les classes minoritaires (stades I et IV).
  \end{itemize}

  \item \textbf{Comparaison des Modèles}
  \begin{itemize}
    \item Tableau comparatif : CNN Baseline vs. EfficientNetB0.
    \item Métriques cibles : Sensibilité $> 95\%$, F2-Score, ROC-AUC par stade.
    \item Sauvegarde du meilleur modèle pour le déploiement.
  \end{itemize}

\end{enumerate}

% -------------------------------------------------------------------
\newpage
\subsection{Jour 4 — Déploiement : FastAPI, Streamlit et Docker}

\begin{daybox}
\textbf{Objectif :} Exposer le modèle via une API robuste, une interface utilisateur et un conteneur Docker prêt à déployer.
\end{daybox}

\vspace{0.4cm}

\begin{enumerate}[label=\textbf{Étape \arabic*.}, leftmargin=2.5cm, itemsep=0.4cm, start=12]

  \item \textbf{API FastAPI — \texttt{api/main.py}}
  \begin{itemize}
    \item Endpoint asynchrone : \texttt{async POST /predict}
    \item Entrée : image médicale (multipart/form-data).
    \item Sortie JSON :
    \begin{itemize}
      \item Stade prédit (classe 0–4)
      \item Score de confiance (probabilité Softmax max)
      \item Heatmap Grad-CAM encodée en base64
      \item Flag \texttt{requires\_review} si confiance $< 0.7$ ou anomalie détectée
    \end{itemize}
    \item Gestion des requêtes simultanées (multi-services hospitaliers) grâce à l'asynchronisme.
  \end{itemize}

  \item \textbf{Dashboard Streamlit — \texttt{dashboard/app.py}}
  \begin{itemize}
    \item Upload d'image médicale (IRM, CT-scan, histopathologie).
    \item Affichage du diagnostic : stade prédit + score de confiance.
    \item Visualisation superposée de la heatmap Grad-CAM.
    \item Indicateur visuel d'alerte : révision manuelle recommandée.
  \end{itemize}

  \item \textbf{Conteneurisation Docker — \texttt{docker/}}
  \begin{itemize}
    \item \texttt{Dockerfile} : image de base \texttt{tensorflow/tensorflow:latest-gpu}.
    \item \texttt{docker-compose.yml} : deux services (\texttt{api} sur le port 8000, \texttt{dashboard} sur le port 8501).
    \item Volumes persistants pour les modèles entraînés.
    \item Garantie de reproductibilité entre environnements dev / production (conforme aux standards des systèmes critiques).
  \end{itemize}

\end{enumerate}

% -------------------------------------------------------------------
\newpage
\section{Considérations Importantes}

\begin{alertbox}
\textbf{\faExclamationTriangle\quad Points d'attention à vérifier avant de commencer}
\end{alertbox}

\vspace{0.5cm}

\begin{enumerate}[label=\textbf{\arabic*.}, leftmargin=1.5cm, itemsep=0.5cm]

  \item \textbf{Téléchargement du Dataset}

  \begin{enumerate}[label=\alph*)]
    \item Créer un compte sur \textbf{kaggle.com} et générer le token API :\\
    \textit{Account $\rightarrow$ API $\rightarrow$ Create New Token} $\Rightarrow$ fichier \texttt{kaggle.json}
    \item Placer le fichier dans \texttt{C:\textbackslash Users\textbackslash <nom>\textbackslash .kaggle\textbackslash}
    \item Télécharger le dataset :
    \begin{verbatim}
pip install kaggle
kaggle datasets download -d masoudnickparvar/brain-tumor-mri-dataset
    \end{verbatim}
    \item Décompresser dans \texttt{data/raw/} puis exécuter \texttt{src/remap\_dataset.py}.
  \end{enumerate}

  \item \textbf{Environnement de développement}

  Le projet IntelliJ actuel est en Java. Il est recommandé de :
  \begin{itemize}
    \item Créer un projet Python dédié (PyCharm ou VS Code) dans le même répertoire.
    \item Utiliser un environnement virtuel : \texttt{python -m venv venv}.
    \item Installer les dépendances : \texttt{pip install -r requirements.txt}.
  \end{itemize}

  \item \textbf{GPU disponible ?}

  \begin{itemize}
    \item \textbf{GPU local disponible} : utiliser \texttt{tensorflow-gpu} + CUDA.
    \item \textbf{Pas de GPU} : entraîner sur \textbf{Google Colab} (gratuit, GPU T4), puis télécharger le modèle \texttt{.h5} pour le déploiement local via FastAPI et Docker.
  \end{itemize}

\end{enumerate}

% -------------------------------------------------------------------
\section{Récapitulatif des Livrables}

\begin{center}
\begin{tabularx}{\textwidth}{|c|l|X|c|}
\hline
\rowcolor{primaryBlue!20}
\textbf{Jour} & \textbf{Fichier} & \textbf{Description} & \textbf{Priorité} \\
\hline
1 & \texttt{src/remap\_dataset.py}   & Remapping 4 classes $\rightarrow$ 5 stades WHO  & \textcolor{dangerRed}{\textbf{Critique}} \\
1 & \texttt{src/data\_pipeline.py}   & Prétraitement, augmentation, encodage      & \textcolor{dangerRed}{\textbf{Critique}} \\
1 & \texttt{src/model\_baseline.py}  & Architecture CNN Sequential Keras          & \textcolor{dangerRed}{\textbf{Critique}} \\
1 & \texttt{src/train.py}            & Entraînement + sauvegarde                  & \textcolor{dangerRed}{\textbf{Critique}} \\
1 & \texttt{src/evaluate.py}         & Métriques cliniques + matrice de confusion & \textcolor{dangerRed}{\textbf{Critique}} \\
\hline
2 & \texttt{src/gradcam.py}          & Heatmaps XAI (Grad-CAM)                   & \textcolor{warningOrange}{\textbf{Élevée}} \\
2 & \texttt{src/autoencoder.py}      & Détection d'anomalies par reconstruction  & \textcolor{warningOrange}{\textbf{Élevée}} \\
\hline
3 & \texttt{src/model\_advanced.py}  & EfficientNetB0 + Fine-Tuning               & \textcolor{accentGreen}{\textbf{Haute}} \\
3 & \texttt{src/losses.py}           & Focal Loss / Weighted Cross-Entropy        & \textcolor{accentGreen}{\textbf{Haute}} \\
\hline
4 & \texttt{api/main.py}             & API FastAPI async                          & \textcolor{accentGreen}{\textbf{Haute}} \\
4 & \texttt{dashboard/app.py}        & Dashboard Streamlit + Grad-CAM             & \textcolor{accentGreen}{\textbf{Haute}} \\
4 & \texttt{docker/Dockerfile}       & Conteneurisation TensorFlow-GPU            & \textcolor{accentGreen}{\textbf{Haute}} \\
\hline
\end{tabularx}
\end{center}

\vspace{1cm}
\begin{center}
\textcolor{primaryBlue}{\large\textbf{Bonne chance — Inch'Allah le projet sera terminé dans les délais ! \faRocket}}
\end{center}

\end{document}

